% ============================================================
%  Chương I – Giới thiệu
%  ~300 từ · 1 trang
% ============================================================
\chapter{Giới thiệu}
\label{chap:intro}

% Ý chính: đặt câu hỏi trung tâm — mật mã bất đối xứng giải quyết được
% key exchange nhưng chưa giải quyết identity binding. PKI là câu trả lời.

% ----------------------------------------------------------
\section{Vấn đề đặt ra}
\label{sec:1.1}

Mật mã khóa công khai (public-key cryptography) đã mang lại một bước tiến quan trọng trong lĩnh vực an toàn thông tin khi cho phép các thực thể trao đổi dữ liệu bí mật và thực hiện chữ ký số mà không cần phải chia sẻ khóa bí mật từ trước. Nhờ đó, các hệ thống có thể hoạt động trên môi trường mạng mở như Internet mà vẫn đảm bảo được tính bảo mật và toàn vẹn của dữ liệu.

Tuy nhiên, một vấn đề cốt lõi vẫn chưa được giải quyết bởi bản thân mật mã khóa công khai, đó là: làm thế nào để xác định rằng một khóa công khai thực sự thuộc về đúng thực thể mà nó tuyên bố đại diện. Nói cách khác, hệ thống vẫn thiếu một cơ chế để ràng buộc giữa danh tính (identity) và khóa công khai (public key).

Trong môi trường mạng không tin cậy, kẻ tấn công có thể khai thác điểm yếu này để thực hiện tấn công trung gian (Man-in-the-Middle). Ví dụ, khi hai bên A và B trao đổi khóa công khai, kẻ tấn công có thể chặn kết nối và thay thế khóa công khai của B bằng khóa của chính mình. Khi đó, A sẽ mã hóa dữ liệu bằng khóa của kẻ tấn công mà không hề hay biết. Mặc dù các thuật toán mật mã vẫn hoạt động chính xác về mặt kỹ thuật, toàn bộ thông tin trao đổi đã bị lộ và có thể bị chỉnh sửa trước khi đến đích.

Một ví dụ thực tế có thể thấy trong quá trình truy cập website. Nếu không có cơ chế xác thực đáng tin cậy, người dùng có thể bị chuyển hướng đến một website giả mạo có giao diện giống hệt website thật. Website giả này cung cấp một khóa công khai do kẻ tấn công kiểm soát, khiến người dùng vô tình thiết lập kết nối “an toàn” với chính kẻ tấn công.

Do đó, một yêu cầu quan trọng được đặt ra là cần có một cơ chế đáng tin cậy để xác minh rằng một khóa công khai thực sự thuộc về đúng thực thể tương ứng. Cơ chế này phải hoạt động hiệu quả ngay cả trong môi trường mạng mở, nơi không thể tin tưởng các kênh truyền thông. Hạ tầng khóa công khai (Public Key Infrastructure – PKI) được đề xuất nhằm giải quyết chính vấn đề này, bằng cách cung cấp một hệ thống quản lý và xác thực khóa công khai dựa trên các bên thứ ba tin cậy.
�� 1.2 Phạm vi và bố cục báo cáo
\section{Phạm vi và bố cục báo cáo}
\label{sec:1.2}

Báo cáo này tập trung nghiên cứu về Hạ tầng khóa công khai (Public Key Infrastructure – PKI), với mục tiêu làm rõ vai trò của PKI trong việc đảm bảo an toàn thông tin trong các hệ thống hiện đại. Nội dung báo cáo tiếp cận từ góc độ hệ thống và ứng dụng, không đi sâu vào các chứng minh toán học phức tạp của các thuật toán mật mã, mà tập trung vào việc giải thích nguyên lý hoạt động, kiến trúc và cách thức triển khai trong thực tế.

Cách tiếp cận của báo cáo được tổ chức theo hướng mỗi chương trả lời một câu hỏi cụ thể, giúp người đọc từng bước xây dựng hiểu biết một cách logic và có hệ thống.

Chương~\ref{chap:crypto} trả lời câu hỏi “Tại sao cần PKI?” thông qua việc trình bày các nền tảng của mật mã đối xứng và bất đối xứng, cũng như các hạn chế còn tồn tại của chúng. 

Chương~\ref{chap:pki-arch} tập trung vào câu hỏi “PKI gồm những thành phần nào?”, trong đó mô tả kiến trúc hệ thống và các thực thể như Certification Authority (CA), Registration Authority (RA) và các cơ chế quản lý chứng chỉ.

Chương~\ref{chap:x509} trả lời câu hỏi “Chứng chỉ số được định nghĩa như thế nào?”, trình bày cấu trúc và vai trò của chứng chỉ số X.509 trong việc ràng buộc giữa danh tính và khóa công khai.

Chương~\ref{chap:deploy} làm rõ “PKI được sử dụng ở đâu trong thực tế?”, với các ví dụ tiêu biểu như HTTPS, chữ ký số và các hệ thống bảo mật khác.

Chương~\ref{chap:opensource} khảo sát các công cụ và nền tảng mã nguồn mở cho phép triển khai PKI trong thực tế, trong khi Chương~\ref{chap:challenges} thảo luận về các thách thức hiện tại cũng như xu hướng phát triển trong tương lai của PKI.
